\documentclass[memoria_final]{subfiles}
\begin{document}
% Deteccion standalone vs. incluido en memoria_final
% \jobname = 'memoria_fase2' cuando se compila solo
% \jobname = 'memoria_final' cuando se incluye desde el documento final
\edef\@currjob{\jobname}
\def\@mainjob{memoria_final}
\ifx\@currjob\@mainjob
  % ---- incluido en memoria_final: solo separador de parte ----
  \clearpage
  \part{Fase 2: ALU, Teclado y Displays}
\else
  % ---- standalone: mostrar portada e indice propios ----
\hypertarget{inicio}{}

% -------------------------------------------------------------------
%  PORTADA
% -------------------------------------------------------------------
\begin{titlepage}
  \thispagestyle{empty}
  \begin{tikzpicture}[remember picture,overlay]
    \fill[white](current page.south west)rectangle(current page.north east);
    \fill[MainBlue]
      (current page.north west)rectangle
      ([yshift=-0.22\paperheight]current page.north east);
    \draw[MainBlue!40,line width=2pt]
      ([yshift=-0.22\paperheight]current page.north west)--
      ([yshift=-0.22\paperheight]current page.north east);
  \end{tikzpicture}

  \vspace*{-1cm}
  \begin{flushleft}
    {\fontsize{9}{12}\selectfont\color{white!80}
      GRADO EN INGENIERÍA ELECTRÓNICA DE COMUNICACIONES}\\[1ex]
    {\fontsize{24}{28}\selectfont\bfseries\color{white}
      Diseño Digital 2}\\[1ex]
    {\fontsize{11}{14}\selectfont\color{white!80}
      Curso académico 2025--2026}
  \end{flushleft}

  \vspace*{0.11\paperheight}
  \begin{center}
    {\Large\color{MainBlue!80}\textbf{BLOQUE II}
      \quad\textbullet\quad ALU, Teclado y Displays}\\[0.8cm]

    {\fontsize{30}{38}\selectfont\bfseries\color{MainBlue}
      Memoria del Proyecto}\\[0.2cm]
    {\fontsize{22}{28}\selectfont\bfseries\color{MainBlue!80}
      Bloque II --- Fase 2}\\[0.5cm]

    {\fontsize{12}{18}\selectfont\color{black!60}
      Diseño, implementación y verificación en VHDL-93}\\[0.2cm]
    {\fontsize{11}{14}\selectfont\ttfamily\color{black!50}
      alu \quad clk\_div \quad timer \quad ctrl\_tec
      \quad interfaz\_teclado \quad displays}\\[0.4cm]

    {\fontsize{11}{14}\selectfont\bfseries\color{MainBlue}
      Versión 1.0 --- Mayo 2026}

    \vspace{1cm}
    \renewcommand{\arraystretch}{1.5}
    \begin{tabular}{@{}r@{\hspace{1em}}l@{}}
      {\small\color{MainBlue}\faLayerGroup}\enspace\textbf{Fase:}
        & Bloque II --- Fase 2 \\
      {\small\color{MainBlue}\faCalendar}\enspace\textbf{Entrega:}
        & 14 de Mayo de 2026 \\
      {\small\color{MainBlue}\faMicrochip}\enspace\textbf{Plataforma:}
        & DECA MAX10 + conector XDECA \\
      {\small\color{MainBlue}\faClock}\enspace\textbf{Reloj:}
        & 50\,MHz \\
    \end{tabular}

    \vspace{1cm}
    \begin{tcolorbox}[
      enhanced,width=0.80\textwidth,
      colback=SoftBlue,colframe=MainBlue,
      boxrule=0.8pt,arc=4pt,
      title={\centering\textbf{\faUsers\quad Miembros del Grupo}},
      coltitle=white,fonttitle=\large]
      \begin{center}
        \renewcommand{\arraystretch}{1.3}
        \begin{tabular}{@{}ll@{}}
          \textbf{Lucía Llana Carmona}   &\small\texttt{lucia.llana@alumnos.upm.es}\\
          \textbf{María Moya Marcilla}    &\small\texttt{maria.moya.marcilla@alumnos.upm.es}\\
          \textbf{Pedro Gil Bolaños}      &\small\texttt{pedro.gilb@alumnos.upm.es}\\
          \textbf{Pedro Osma Moya}        &\small\texttt{p.osma@alumnos.upm.es}\\
          \textbf{Eduardo Muñoz de Maya}  &\small\texttt{eduardo.mdemaya@alumnos.upm.es}\\
        \end{tabular}
      \end{center}
    \end{tcolorbox}
  \end{center}
  \vfill
  \begin{center}
    {\fontsize{8}{10}\selectfont\color{black!40}
      Grado en Ingeniería Electrónica de Comunicaciones
      \quad$\cdot$\quad Diseño Digital 2}
  \end{center}
  \vspace{0.5cm}
\end{titlepage}
\newpage
\tableofcontents
\newpage
\fi% fin \ifx\@currjob\@mainjob
\fancyhead[R]{\small\textbf{Memoria --- Fase 2 (ALU, Teclado, Displays)}}

% ===================================================================
\section{Introducción y objetivos}
% ===================================================================

En esta segunda fase se han diseñado, implementado y verificado los bloques
intermedios de la calculadora: la \textbf{unidad aritmético-lógica} (ALU),
el \textbf{sistema de temporización}, el \textbf{controlador de teclado
matricial} y el \textbf{controlador de displays multiplexados}. Estos módulos
constituyen el núcleo de la \textit{capa de interfaz} del sistema.

El objetivo de la fase es disponer de un camino de datos completamente funcional
y simulable que pueda integrarse directamente con el controlador (FSM) y los
conversores de código desarrollados en la Fase~1.

\vspace{0.5em}
\begin{tcolorbox}[colback=SoftBlue,colframe=MainBlue,boxrule=0.8pt,arc=2pt,
  left=0.5cm,right=0.5cm,top=0.3cm,bottom=0.3cm]
  \small\faInfoCircle\enspace\textbf{Criterio de codificación adoptado en todo
  el proyecto:}
  todos los módulos RTL emplean \textbf{señales} en lugar de variables.
  Una señal modela directamente un cable o un registro físico: su valor solo se
  actualiza al término del proceso (equivalente al flanco de reloj), igual que
  ocurre en el circuito real. Por el contrario, una variable se actualiza de
  forma inmediata dentro del proceso---comportamiento propio del software---lo
  que dificulta la correspondencia con el hardware y oculta la señal al
  simulador (no aparece en las formas de onda). Este criterio se justifica con
  detalle en el Apéndice~\ref{sec:senales-variables}.
\end{tcolorbox}

% ===================================================================
\section{Arquitectura global del sistema}
% ===================================================================

La figura~\ref{fig:arq-global} muestra la jerarquía completa con todos los
módulos, señales de interconexión y agrupación por fases de desarrollo.

\begin{figure}[H]
\centering
\begin{tikzpicture}[node distance=0.65cm and 1.0cm,font=\sffamily\footnotesize]

  % ---- Fila 1: entrada -------------------------------------------------
  \node[bloqueg,text width=2.2cm](hwtec)
    {Teclado\\4$\times$4\\{\tiny XDECA}};

  \node[bloqueg,text width=2.6cm,right=1.0cm of hwtec](itec)
    {\texttt{interfaz\_}\\\texttt{teclado}\\[2pt]
     {\tiny clk\_div + ctrl\_tec}};

  \node[bloqueo,text width=2.6cm,right=1.2cm of itec](ctrl)
    {\texttt{controlador}\\[2pt]
     {\tiny FSM principal\\Fase~3}};

  % ---- Fila 2: conversores BCD->bin ------------------------------------
  \node[bloque,text width=2.3cm,
        below right=1.0cm and 0.2cm of ctrl](b2b1)
    {\texttt{bcd\_to\_bin}\\[2pt]{\tiny op1}};

  \node[bloque,text width=2.3cm,
        right=0.5cm of b2b1](b2b2)
    {\texttt{bcd\_to\_bin}\\[2pt]{\tiny op2}};

  % ---- Fila 3: ALU -----------------------------------------------------
  \node[bloque,text width=2.6cm,
        below=1.1cm of $(b2b1)!0.5!(b2b2)$](alu)
    {\texttt{alu}\\[2pt]
     {\tiny suma / resta / mult\\comp.\ a~2}};

  \node[bloqueip,below=0.45cm of alu](mult)
    {\texttt{lpm\_mult}\\{\tiny IP Altera\\$11{\times}11$\,b}};

  % ---- Fila 4: bin_to_bcd + timer -------------------------------------
  \node[bloque,text width=2.4cm,
        below left=1.5cm and 1.0cm of alu](b2bcd)
    {\texttt{bin\_to\_bcd}\\[2pt]
     {\tiny Horner BCD\\20\,ciclos}};

  \node[bloque,text width=2.2cm,
        right=2.4cm of b2bcd](timer)
    {\texttt{timer}\\[2pt]
     {\tiny tic\_1ms\\tic\_5ms}};

  % ---- Fila 5: displays -----------------------------------------------
  \node[bloqueg,text width=2.4cm,
        below=1.1cm of $(b2bcd)!0.5!(timer)$](disp)
    {\texttt{displays}\\[2pt]
     {\tiny mux 8$\times$7\,seg\\125\,Hz}};

  \node[bloqueg,text width=2.0cm,
        right=1.0cm of disp](hwdisp)
    {Displays\\8$\times$7\,seg\\{\tiny DECA}};

  % ---- Flechas: teclado -----------------------------------------------
  \draw[flecha](hwtec)  -- node[above,font=\tiny]{col/fila} (itec);
  \draw[flecha](itec)   -- node[above,font=\tiny]{tecla[3:0]\\tp} (ctrl);

  % ---- Flechas: datapath ----------------------------------------------
  \draw[flecha](ctrl.south east)
    -- node[right,font=\tiny,xshift=2pt]{op1\_bcd} (b2b1.north);
  \draw[flecha](ctrl.south east)
    -- node[right,font=\tiny,xshift=2pt]{op2\_bcd} (b2b2.north);
  \draw[bus](b2b1.south) -- node[left,font=\tiny]{10\,b} (alu.north west);
  \draw[bus](b2b2.south) -- node[right,font=\tiny]{10\,b} (alu.north east);
  \draw[bus](alu.south)  -- node[right,font=\tiny]{20\,b} (b2bcd.north east);
  \draw[bus](b2bcd.south east) -- node[above,font=\tiny]{BCD\,24\,b} (disp.north west);
  \draw[flecha](timer.south) -- node[right,font=\tiny]{tic\_1ms} (disp.north east);
  \draw[flecha](disp) -- (hwdisp);

  % ---- Flechas: control -----------------------------------------------
  \draw[flechaO](ctrl.south)
    |- node[right,font=\tiny,yshift=4pt]{start\_bcd} (b2bcd.west);
  \draw[flechaO](b2bcd.north)
    -- node[left,font=\tiny]{done} ++(0,0.5)
    -| (ctrl.south);
  \draw[flechaO](ctrl.south)
    |- node[below,font=\tiny]{pres / res\_sgn} (disp.west);
  \draw[flechaO](alu.west)
    -- node[above,font=\tiny]{op\_sel} ++(-0.6,0)
    |- (ctrl.west);

  % ---- ALU <-> lpm_mult -----------------------------------------------
  \draw[flecha,color=gray](alu.south) -- (mult.north);
  \draw[flecha,color=gray](mult.north) -- (alu.south);

  % ---- Marcos de fase -------------------------------------------------
  \begin{scope}[on background layer]
    \node[nucleo,fit=(b2b1)(b2b2)(b2bcd),
      label={[font=\tiny\bfseries\color{MainBlue},yshift=2pt]above:Fase~1}]{};
    \node[nucleo,fit=(itec)(alu)(timer)(disp),
      label={[font=\tiny\bfseries\color{MainBlue},yshift=2pt]above:Fase~2}]{};
    \node[nucleo,fit=(ctrl),
      label={[font=\tiny\bfseries\color{MainOrange},yshift=2pt]above:Fase~3}]{};
  \end{scope}

\end{tikzpicture}
\caption{Arquitectura jerárquica completa de la calculadora.}
\label{fig:arq-global}
\end{figure}
\FloatBarrier

\subsection{Señales de interconexión principales}

\begin{table}[H]
  \centering
  \caption{Señales clave entre módulos}
  \begin{tabular}{llll}
    \toprule
    \textbf{Señal} & \textbf{Ancho} & \textbf{De} & \textbf{A} \\
    \midrule
    \texttt{tecla[3:0]}       & 4\,b  & \texttt{ctrl\_tec}    & \texttt{controlador} \\
    \texttt{tecla\_pulsada}   & 1\,b  & \texttt{ctrl\_tec}    & \texttt{controlador} \\
    \texttt{op1\_bcd[11:0]}   & 12\,b & \texttt{controlador}  & \texttt{bcd\_to\_bin}, \texttt{displays} \\
    \texttt{op2\_bcd[11:0]}   & 12\,b & \texttt{controlador}  & \texttt{bcd\_to\_bin}, \texttt{displays} \\
    \texttt{op1\_bin[10:0]}   & 11\,b & lógica inline C2      & \texttt{alu} \\
    \texttt{op2\_bin[10:0]}   & 11\,b & lógica inline C2      & \texttt{alu} \\
    \texttt{res[19:0]}        & 20\,b & \texttt{alu}          & \texttt{bin\_to\_bcd} \\
    \texttt{res\_sgn}         & 1\,b  & \texttt{alu}          & \texttt{controlador}, \texttt{displays} \\
    \texttt{bcd[23:0]}        & 24\,b & \texttt{bin\_to\_bcd} & \texttt{displays} \\
    \texttt{done}             & 1\,b  & \texttt{bin\_to\_bcd} & \texttt{controlador} \\
    \texttt{tic\_1ms}         & 1\,b  & \texttt{timer}        & \texttt{displays} \\
    \texttt{pres[1:0]}        & 2\,b  & \texttt{controlador}  & \texttt{displays} \\
    \bottomrule
  \end{tabular}
\end{table}
\FloatBarrier

% ===================================================================
\section{Unidad Aritmético-Lógica (\texttt{alu.vhd})}
% ===================================================================

\subsection{Especificación}

\begin{tcolorbox}[colback=SoftGreen,colframe=MainGreen,boxrule=0.8pt,arc=2pt,
  left=0.5cm,right=0.5cm,top=0.3cm,bottom=0.3cm]
  \small
  \begin{tabular}{ll}
    \textbf{Tipo:}      & Lógica combinacional (sin reloj).\\[3pt]
    \textbf{Entradas:}  &
      \texttt{op1[10:0]}, \texttt{op2[10:0]} en complemento a 2,
      rango $[-999,\,+999]$;\\
    & \texttt{op[1:0]}: ``\texttt{00}''=suma, ``\texttt{01}''=resta,
      ``\texttt{10}''=multiplicación.\\[3pt]
    \textbf{Salidas:}   &
      \texttt{res[19:0]} magnitud del resultado;
      \texttt{res\_sgn} signo (\texttt{'1'} = negativo);\\
    & \texttt{overflow} (fijado a \texttt{'0'}: los operandos $\pm999$
      no pueden desbordar 20\,bits).\\[3pt]
    \textbf{Rango máx.} & $999\times999 = 998\,001 < 2^{20}$.
  \end{tabular}
\end{tcolorbox}

\subsection{Camino de datos}

La figura siguiente muestra el camino de datos interno de la ALU.

\begin{center}
\begin{tikzpicture}[node distance=0.6cm and 1.3cm,font=\sffamily\footnotesize]

  % Entradas
  \node[bloque,text width=1.8cm](op1n){op1[10:0]\\{\tiny C2, 11\,b}};
  \node[bloque,text width=1.8cm,right=0.3cm of op1n](op2n){op2[10:0]\\{\tiny C2, 11\,b}};

  % Extensión de signo
  \node[bloqueo,text width=2.8cm,below=0.8cm of $(op1n)!0.5!(op2n)$](ext)
    {Extensión de signo\\{\tiny 11\,b $\to$ 12\,b}\\
     {\tiny \texttt{op1(10)\&op1}}};

  % Operaciones
  \node[bloque,text width=1.6cm,below left=0.8cm and 0.5cm of ext](suma)
    {Suma\\{\tiny \texttt{op1\_12 + op2\_12}}\\12\,b};
  \node[bloque,text width=1.6cm,right=0.3cm of suma](resta)
    {Resta\\{\tiny \texttt{op1\_12 - op2\_12}}\\12\,b};
  \node[bloqueip,text width=1.6cm,right=0.3cm of resta](mult)
    {\texttt{lpm\_mult}\\{\tiny $11\times11 \to 22$\,b}};

  % Complementos a 2
  \node[bloqueg,text width=4.5cm,below=1.0cm of $(suma)!0.5!(resta)$](c2)
    {Extracción de magnitud\\{\tiny
     \texttt{sum\_neg} $\leftarrow$ \texttt{(not sum\_12)+1}\\
     \texttt{sub\_neg} $\leftarrow$ \texttt{(not sub\_12)+1}\\
     \texttt{mul\_neg} $\leftarrow$ \texttt{(not mult\_res)+1}}};

  % MUX de salida
  \node[bloqueo,text width=3.8cm,below=0.8cm of c2](mux)
    {Selección \texttt{when/else}\\{\tiny según \texttt{op} y bit de signo}};

  % Salidas
  \node[bloque,text width=1.6cm,below left=0.8cm and 0.3cm of mux](sres)
    {\texttt{res[19:0]}\\{\tiny magnitud}};
  \node[bloque,text width=1.6cm,right=0.3cm of sres](ssgn)
    {\texttt{res\_sgn}\\{\tiny signo}};

  % Flechas
  \draw[flecha](op1n.south)--(ext.north west);
  \draw[flecha](op2n.south)--(ext.north east);
  \draw[flecha](ext.south west)--(suma.north);
  \draw[flecha](ext.south)--(resta.north);
  \draw[flecha](op1n.south east)--(mult.north west);
  \draw[flecha](op2n.south east)--(mult.north);
  \draw[flecha](suma.south)--(c2.north west);
  \draw[flecha](resta.south)--(c2.north);
  \draw[flecha](mult.south)--(c2.north east);
  \draw[flecha](c2.south)--(mux.north);
  \draw[flecha](mux.south west)--(sres.north);
  \draw[flecha](mux.south east)--(ssgn.north);
\end{tikzpicture}
\end{center}
\FloatBarrier

\subsection{Implementación en VHDL}

Los operandos se extienden de 11 a 12 bits replicando el bit de signo (MSB),
lo que permite realizar suma y resta directamente en aritmética de complemento
a 2 con un bit extra que actúa como indicador de signo del resultado:

\begin{lstlisting}[caption={Extensión de signo y operaciones aritméticas en \texttt{alu.vhd}}]
-- Extension de signo: 11 -> 12 bits (replica op(10) en la posicion 11)
op1_12 <= op1(10) & op1;
op2_12 <= op2(10) & op2;

-- Suma y resta en complemento a 2 (12 bits)
sum_12 <= op1_12 + op2_12;
sub_12 <= op1_12 - op2_12;
\end{lstlisting}

Para la multiplicación se instancia el componente \texttt{lpm\_mult}, que es
el modelo de simulación de una IP de Altera (multiplicador con signo de
$11 \times 11 \to 22$ bits):

\begin{lstlisting}[caption={Instanciación del multiplicador IP en \texttt{alu.vhd}}]
U_MULT: lpm_mult
    port map(
        dataa  => op1,
        datab  => op2,
        result => mult_res    -- 22 bits, complemento a 2
    );
\end{lstlisting}

La extracción de la magnitud y el signo se realiza íntegramente con
\textbf{asignaciones concurrentes} \texttt{when/else}, sin ningún proceso ni
variable. Primero se calculan los complementos a~2 de los resultados posibles
(para tener la magnitud en el caso negativo), y luego se selecciona:

\begin{lstlisting}[caption={Extracción de magnitud sin variables ni proceso en \texttt{alu.vhd}}]
-- Complementos a 2: magnitud cuando el resultado es negativo
sum_neg <= (not sum_12)   + 1;
sub_neg <= (not sub_12)   + 1;
mul_neg <= (not mult_res) + 1;

-- Seleccion de magnitud segun operacion y bit de signo del resultado
res <= "00000000" & sum_neg     when op="00" and sum_12(11)  ='1' else
       "00000000" & sum_12      when op="00"                      else
       "00000000" & sub_neg     when op="01" and sub_12(11)  ='1' else
       "00000000" & sub_12      when op="01"                      else
       mul_neg(19 downto 0)     when op="10" and mult_res(21)='1' else
       mult_res(19 downto 0);

-- Signo: '1' cuando el resultado es negativo
res_sgn <= '1' when (op="00" and sum_12(11)  ='1') or
                    (op="01" and sub_12(11)  ='1') or
                    (op="10" and mult_res(21)='1') else '0';
\end{lstlisting}

\begin{tcolorbox}[colback=SoftOrange,colframe=MainOrange,boxrule=0.8pt,arc=2pt,
  left=0.5cm,right=0.5cm,top=0.3cm,bottom=0.3cm]
  \small\faLightbulb\enspace\textbf{¿Por qué \texttt{when/else} y no un
  proceso?} Porque la ALU es combinacional: no hay reloj ni estado. Un
  proceso con lista de sensibilidad sería equivalente, pero las asignaciones
  concurrentes expresan directamente la intención --- ``este cable vale X o
  Y dependiendo de la condición'' --- sin necesitar variables intermedias.
\end{tcolorbox}

\subsection{Verificación: \texttt{tb\_alu.vhd}}

El testbench es puramente combinacional (sin reloj). Cambia las entradas cada
20\,ns y comprueba el resultado mediante la función auxiliar \texttt{check},
que genera un \texttt{report} de error si la magnitud o el signo no coinciden
con los esperados:

\begin{lstlisting}[caption={Procedimiento \texttt{check} del testbench \texttt{tb\_alu}}]
procedure check(
    tag     : in string;
    got_mag : in std_logic_vector(19 downto 0);
    got_sgn : in std_logic;
    exp_mag : in integer;
    exp_sgn : in std_logic
) is
begin
    assert got_sgn = exp_sgn
        report tag & ": signo incorrecto." severity error;
    assert got_mag = exp_mag
        report tag & ": magnitud incorrecta." severity error;
end procedure;
\end{lstlisting}

\begin{table}[H]
  \centering
  \caption{Casos de prueba del testbench \texttt{tb\_alu} (10 casos)}
  \begin{tabular}{rrlrcc}
    \toprule
    \textbf{op1} & \textbf{op2} & \textbf{Operación} &
    \textbf{Magnitud esperada} & \textbf{Signo} &
    \textbf{Propósito} \\
    \midrule
    $+100$ & $+200$ & suma  & 300     & $+$ & Suma positivos \\
    $-100$ & $-200$ & suma  & 300     & $-$ & Suma negativos \\
    $-100$ & $+200$ & suma  & 100     & $+$ & Suma signos distintos \\
    $+300$ & $+100$ & resta & 200     & $+$ & Resta normal \\
    $+100$ & $+300$ & resta & 200     & $-$ & Resultado negativo \\
    $-100$ & $+200$ & resta & 300     & $-$ & Resta con negativo \\
    $+12$  & $+12$  & mult  & 144     & $+$ & Producto simple \\
    $-3$   & $+5$   & mult  & 15      & $-$ & Producto negativo \\
    $+999$ & $+999$ & mult  & 998\,001 & $+$ & Valor máximo posible \\
    $0$    & $+999$ & mult  & 0       & $+$ & Cero en multiplicación \\
    \bottomrule
  \end{tabular}
\end{table}

\noindent La simulación termina con \textbf{Errors: 0, Warnings: 0} en
ModelSim~10.4d con VHDL-93.
\FloatBarrier

% ===================================================================
\section{Temporización: \texttt{clk\_div.vhd} y \texttt{timer.vhd}}
% ===================================================================

\subsection{Divisor de reloj genérico (\texttt{clk\_div})}

Genera un pulso \texttt{tic} de \textbf{exactamente un ciclo de reloj} de
anchura cada \texttt{DIV+1} ciclos. El genérico \texttt{DIV} permite
instanciar el mismo módulo para distintos periodos:

\begin{lstlisting}[caption={Implementación de \texttt{clk\_div.vhd}}]
process(clk, nRst)
begin
    if nRst = '0' then
        cnt     <= 0;
        tic_reg <= '0';
    elsif clk'event and clk = '1' then
        tic_reg <= '0';          -- por defecto pulso bajo
        if cnt = DIV then
            tic_reg <= '1';      -- pulso de 1 ciclo
            cnt     <= 0;
        else
            cnt <= cnt + 1;
        end if;
    end if;
end process;
tic <= tic_reg;                  -- asignacion concurrente fuera del proceso
\end{lstlisting}

\begin{table}[H]
  \centering
  \caption{Instancias de \texttt{clk\_div} en el proyecto}
  \begin{tabular}{lcc}
    \toprule
    \textbf{Uso} & \textbf{DIV} & \textbf{Periodo resultante} \\
    \midrule
    Tic teclado (\texttt{interfaz\_teclado}) & 249\,999 & 5\,ms a 50\,MHz \\
    Simulación (genérico reducido)           & 4        & 100\,ns \\
    \bottomrule
  \end{tabular}
\end{table}

\subsection{Temporizador doble (\texttt{timer})}

Genera dos pulsos independientes a partir del mismo reloj de 50\,MHz:
\texttt{tic\_1ms} para la multiplexación de displays y \texttt{tic\_5ms}
como referencia temporal adicional.

\begin{lstlisting}[caption={Generación de \texttt{tic\_1ms} y \texttt{tic\_5ms} en \texttt{timer.vhd}}]
-- Tic cada 1 ms
if cnt_1ms = DIV_1ms then
    t1ms_reg <= '1';
    cnt_1ms  <= 0;
    -- Tic cada 5 ms (cada 5 tics de 1 ms)
    if cnt_5ms = DIV_5ms then
        t5ms_reg <= '1';
        cnt_5ms  <= 0;
    else
        cnt_5ms <= cnt_5ms + 1;
    end if;
else
    cnt_1ms <= cnt_1ms + 1;
end if;
\end{lstlisting}

Los dos contadores son completamente independientes en el proceso y se resetean
con el mismo \texttt{nRst} activo bajo.

% ===================================================================
\section{Controlador de teclado (\texttt{ctrl\_tec.vhd})}
% ===================================================================

\subsection{Descripción general y especificación de puertos}

El módulo gestiona un teclado matricial 4$\times$4 conectado al conector XDECA.
La interfaz utiliza buses de 4 bits en lugar de señales individuales, lo que
simplifica el cableado en la jerarquía superior:

\begin{tcolorbox}[colback=SoftGreen,colframe=MainGreen,boxrule=0.8pt,arc=2pt,
  left=0.5cm,right=0.5cm,top=0.3cm,bottom=0.3cm]
  \small
  \begin{tabular}{lll}
    \textbf{Puerto} & \textbf{Dirección} & \textbf{Descripción} \\
    \hline
    \texttt{clk}           & entrada & Reloj 50\,MHz \\
    \texttt{nRst}          & entrada & Reset activo bajo \\
    \texttt{tic}           & entrada & Pulso 5\,ms (de \texttt{clk\_div}) \\
    \texttt{columna[3:0]}  & entrada & Columnas activo bajo \\
    \texttt{fila[3:0]}     & buffer  & Filas de escaneo, activo bajo \\
    \texttt{tecla[3:0]}    & buffer  & Código hex de la última tecla \\
    \texttt{tecla\_pulsada}& buffer  & Pulso 1 ciclo al soltar la tecla \\
    \texttt{pulso\_largo}  & buffer  & Activo tras 2\,s mantenida \\
  \end{tabular}\\[4pt]
  \small\textbf{Genérico:} \texttt{TICS\_2s = 400}
  (400 tics $\times$ 5\,ms = 2\,s a 50\,MHz).
\end{tcolorbox}

\subsection{Mapa del teclado}

La distribución física del teclado 4$\times$4 se codifica en complemento
one-hot activo bajo para filas y columnas:

\begin{table}[H]
  \centering
  \caption{Distribución de teclas: código hex asignado a cada posición}
  \begin{tabular}{ccccc}
    \toprule
    \textbf{fila} & \textbf{col\,``1110''} & \textbf{col\,``1101''} &
    \textbf{col\,``1011''} & \textbf{col\,``0111''} \\
    \midrule
    \texttt{"1110"} fila~0 & \texttt{1} & \texttt{2} & \texttt{3} & \texttt{F} \\
    \texttt{"1101"} fila~1 & \texttt{4} & \texttt{5} & \texttt{6} & \texttt{E}~(×) \\
    \texttt{"1011"} fila~2 & \texttt{7} & \texttt{8} & \texttt{9} & \texttt{D}~(−) \\
    \texttt{"0111"} fila~3 & \texttt{A}~(+) & \texttt{0} & \texttt{B}~(=) & \texttt{C}~(±) \\
    \bottomrule
  \end{tabular}
\end{table}

\subsection{Mecanismo de escaneo de filas}

El módulo activa una fila a la vez poniendo su bit a \texttt{'0'} (activo bajo),
manteniendo las demás a \texttt{'1'}. En cada tic de 5\,ms rota el patrón
one-hot hacia la izquierda:

$$\texttt{fila} \;:\; \texttt{"1110"} \to \texttt{"1101"} \to \texttt{"1011"} \to \texttt{"0111"} \to \texttt{"1110"} \to \cdots$$

\begin{lstlisting}[caption={Rotación one-hot de filas en \texttt{ctrl\_tec.vhd}}]
process(clk, nRst)
begin
    if nRst = '0' then
        fila <= (others => '1');   -- todas inactivas en reset
    elsif clk'event and clk = '1' then
        if tic = '1' then
            if ena_muestreo = '1' then
                if fila = X"F" then
                    fila <= "1110";                      -- inicio del escaneo
                else
                    fila <= fila(2 downto 0) & fila(3);  -- rotacion izquierda
                end if;
            end if;
        end if;
    end if;
end process;
\end{lstlisting}

El escaneo se pausa (\texttt{ena\_muestreo = '0'}) mientras se detecta
actividad en las columnas, para evitar capturar una columna en medio del
cambio de fila.

\subsection{Antirebote integrado}

No se usa ningún módulo externo de antirebote. La columna leída se registra
en dos etapas:

\begin{enumerate}[leftmargin=1.8cm,itemsep=2pt]
  \item \textbf{\texttt{col\_reb}}: muestreo síncrono de \texttt{columna} en
    cada ciclo de reloj (elimina los metaestables por cruce de dominios de
    reloj).
  \item \textbf{\texttt{col\_ok}}: se actualiza con el valor de
    \texttt{col\_reb} solo en cada tic de 5\,ms, garantizando que el valor
    lleva al menos 5\,ms estable antes de ser aceptado.
\end{enumerate}

\begin{lstlisting}[caption={Antirebote de dos etapas en \texttt{ctrl\_tec.vhd}}]
process(clk, nRst)
begin
    if nRst = '0' then
        col_reb <= (others => '1');
        col_ok  <= (others => '1');
    elsif clk'event and clk = '1' then
        col_reb <= columna;           -- etapa 1: sincronizar
        if tic = '1' then
            col_ok <= col_reb;        -- etapa 2: validar cada 5 ms
        end if;
    end if;
end process;
\end{lstlisting}

\subsection{Detección de pulsación corta y larga}

La señal interna \texttt{ctrl\_tecla\_pulsada} indica si hay alguna columna
activa en \texttt{col\_ok}. La lógica de detección es la siguiente:

\begin{itemize}[leftmargin=1.6cm,itemsep=3pt]
  \item \textbf{Pulsación corta} (\texttt{tecla\_pulsada}): pulso de
    1~ciclo generado en el \textit{flanco de bajada} de
    \texttt{ctrl\_tecla\_pulsada} (es decir, cuando se suelta la tecla),
    siempre que no haya habido pulsación larga:
    \begin{lstlisting}[numbers=none,aboveskip=4pt,belowskip=4pt]
tecla_pulsada <= '1' when pulso_hab_tecla     = '1'
                       and ctrl_tecla_pulsada = '0'
                       and pulso_hab_largo    = '0'
                 else '0';
    \end{lstlisting}
  \item \textbf{Pulsación larga} (\texttt{pulso\_largo}): se activa cuando
    la tecla lleva pulsada \texttt{TICS\_2s} = 400 tics $\times$ 5\,ms
    = 2\,s. En ese caso se suprime el evento de \textit{release} para evitar
    registrar una pulsación corta fantasma.
\end{itemize}

\begin{tcolorbox}[colback=SoftBlue,colframe=MainBlue,boxrule=0.8pt,arc=2pt,
  left=0.5cm,right=0.5cm,top=0.3cm,bottom=0.3cm]
  \small\faInfoCircle\enspace\textbf{¿Por qué detectar en el release y no en
  el press?} Detectar en la pulsación produce problemas si el usuario mantiene
  la tecla más de lo necesario (se registran múltiples pulsaciones). Detectar
  en el release garantiza exactamente un evento por pulsación, y permite
  además distinguir pulsaciones cortas de largas antes de generar el evento.
\end{tcolorbox}

\subsection{Decodificador fila × columna}

El decodificador es un registro síncrono con reset que evalúa la combinación
de \texttt{fila} y \texttt{col\_ok} y guarda el código de tecla en
\texttt{tecla\_aux}. El resultado se transfiere a la salida \texttt{tecla}
únicamente en el momento de la pulsación (corta o larga):

\begin{lstlisting}[caption={Decodificador fila/columna (fragmento) en \texttt{ctrl\_tec.vhd}}]
process(clk, nRst)
begin
    if nRst = '0' then
        tecla_aux <= (others => '0');
    elsif clk'event and clk = '1' then
        if    fila="1110" and col_ok="1110" then tecla_aux <= X"1";
        elsif fila="1110" and col_ok="1101" then tecla_aux <= X"2";
        -- ... (16 combinaciones en total)
        elsif fila="0111" and col_ok="0111" then tecla_aux <= X"C";
        end if;
    end if;
end process;
\end{lstlisting}

\subsection{Jerarquía: \texttt{interfaz\_teclado.vhd}}

El módulo \texttt{interfaz\_teclado} agrupa \texttt{clk\_div} y
\texttt{ctrl\_tec} en un único bloque con la interfaz estándar del sistema.
Los puertos de tipo \texttt{buffer} de \texttt{ctrl\_tec} se conectan a través
de señales internas para cumplir la restricción de VHDL-93 (no se puede
conectar un puerto \texttt{buffer} directamente a un puerto \texttt{out} del
módulo padre):

\begin{lstlisting}[caption={Integración en \texttt{interfaz\_teclado.vhd}}]
signal fila_s, tec_s : std_logic_vector(3 downto 0);
signal tp_s, pl_s    : std_logic;

U_DIV:  entity work.clk_div(rtl)
    generic map(DIV => DIV_TEC)
    port map(clk => clk, nRst => nRst, tic => tic_s);

U_CTRL: entity work.ctrl_tec(rtl)
    generic map(TICS_2s => TICS_2s)
    port map(
        clk => clk, nRst => nRst, tic => tic_s,
        columna => columna,     fila  => fila_s,
        tecla_pulsada => tp_s,  tecla => tec_s,
        pulso_largo   => pl_s
    );

-- Señales intermedias -> salidas del modulo padre
fila <= fila_s;  tecla <= tec_s;
tecla_pulsada <= tp_s;  pulso_largo <= pl_s;
\end{lstlisting}

% ===================================================================
\section{Controlador de displays (\texttt{displays.vhd})}
% ===================================================================

\subsection{Especificación}

\begin{tcolorbox}[colback=SoftGreen,colframe=MainGreen,boxrule=0.8pt,arc=2pt,
  left=0.5cm,right=0.5cm,top=0.3cm,bottom=0.3cm]
  \small
  \begin{tabular}{ll}
    \textbf{Hardware:}     &
      8 displays de 7 segmentos, cátodo común, activo alto (DECA + XDECA).\\[3pt]
    \textbf{Multiplexación:} &
      Puerto \texttt{mux\_disp} (buffer) one-hot activo bajo; rota en cada
      \texttt{tic\_1ms}.\\
    & Frecuencia de refresco: $1000\,\text{Hz} / 8 = 125\,\text{Hz}$
      (sin parpadeo visible).\\[3pt]
    \textbf{Formato:} &
      \texttt{disp[7:0]}: bits \textit{dp, a, b, c, d, e, f, g}
      (bit 7 = punto decimal, bit 0 = g).\\[3pt]
    \textbf{Selector \texttt{pres[1:0]}:} &
      \texttt{"00"} = op1 en D4--D1;\quad
      \texttt{"01"} = op2 en D4--D1;\quad
      \texttt{"10"} = resultado en D8--D1.
  \end{tabular}
\end{tcolorbox}

\subsection{Interfaz de puertos}

Los puertos \texttt{mux\_disp} y \texttt{disp} se declaran como \texttt{buffer}
(no \texttt{out}) porque su valor se lee internamente en la lógica de
multiplexación. Los puertos de datos se nombran \texttt{op1}, \texttt{op2}
y \texttt{res} (sin sufijo \texttt{\_bcd}).

\subsection{Multiplexación}

El registro \texttt{mux\_disp} es un one-hot activo bajo que rota hacia la
derecha en cada \texttt{tic\_1ms}, activando un display diferente en cada
ciclo de refresco. Al ser \texttt{buffer}, puede leerse directamente en la
asignación concurrente de \texttt{dig\_activo}:

\begin{lstlisting}[caption={Registro de multiplexación en \texttt{displays.vhd}}]
catodos: process(clk, nRst)
begin
  if nRst = '0' then
    mux_disp <= (0 => '0', others => '1');
  elsif clk'event and clk = '1' then
    if tic_1ms = '1' then
      mux_disp <= mux_disp(6 downto 0) & mux_disp(7);
    end if;
  end if;
end process catodos;
\end{lstlisting}

\subsection{Preparación y supresión de ceros}

Los tres vectores \texttt{op1\_c}, \texttt{op2\_c} y \texttt{res\_c} de
32\,bits empaquetan, con asignaciones concurrentes \texttt{when/else}, tanto
el identificador de modo (nibble alto) como los dígitos a mostrar con
supresión de ceros no significativos. El nibble de signo se forma
concatenando \texttt{"111"\&op\_sgn}, de modo que \texttt{x"E"} representa
blanco y \texttt{x"F"} el guión negativo:

\begin{lstlisting}[caption={Supresión de ceros en \texttt{displays.vhd} (fragmento)}]
sig1 <= "111" & op1_sgn;

op1_c <=
  "0001"&"1110"&"1110"&"1110"&"1110"&"1110" & sig1 & op1(3 downto 0)
    when op1(11 downto 4) = 0 else
  "0001"&"1110"&"1110"&"1110"&"1110" & sig1 & op1(7 downto 4) & op1(3 downto 0)
    when op1(11 downto 8) = 0 else
  "0001"&"1110"&"1110"&"1110" & sig1 & op1(11 downto 8)
                               & op1(7 downto 4) & op1(3 downto 0);
\end{lstlisting}

\subsection{Decodificación 4\,bits $\to$ 7 segmentos}

La decodificación se realiza mediante un \texttt{process} con \texttt{case}
directamente sobre la señal \texttt{dig\_activo}, sin función auxiliar:

\begin{table}[H]
  \centering
  \caption{Caracteres especiales del display}
  \begin{tabular}{clcc}
    \toprule
    \textbf{Código} & \textbf{Carácter} &
    \textbf{Patrón (dp,a..g)} & \textbf{Uso} \\
    \midrule
    \texttt{x"F"} & \texttt{-} (guión) & \texttt{00000001} & Signo negativo \\
    \texttt{x"E"} & blanco             & \texttt{00000000} & Supresión de ceros \\
    \texttt{x"A"} & \texttt{A}         & \texttt{01110111} & Dígito hex A \\
    \texttt{x"B"} & \texttt{P}         & \texttt{01100111} & Reservado \\
    \texttt{x"C"} & \texttt{C}         & \texttt{01001110} & Reservado \\
    \texttt{x"D"} & \texttt{r}         & \texttt{00000101} & Reservado \\
    \bottomrule
  \end{tabular}
\end{table}
\FloatBarrier

% ===================================================================
\section{Resultados de compilación y simulación}
% ===================================================================

Todos los módulos de la Fase~2 compilan correctamente en ModelSim~10.4d
(Altera ASE) con la flag \texttt{-93} (VHDL-93):

\begin{table}[H]
  \centering
  \caption{Resultado de compilación de todos los módulos}
  \begin{tabular}{llcc}
    \toprule
    \textbf{Fichero} & \textbf{Entidad / Arquitectura} &
    \textbf{Errores} & \textbf{Warnings} \\
    \midrule
    \texttt{fase1/bcd\_to\_bin.vhd}      & \texttt{bcd\_to\_bin / rtl}        & 0 & 0 \\
    \texttt{fase1/bin\_to\_bcd.vhd}      & \texttt{bin\_to\_bcd / rtl}        & 0 & 0 \\
    \texttt{fase2/lpm\_mult.vhd}         & \texttt{lpm\_mult / behavioral}    & 0 & 0 \\
    \texttt{fase2/alu.vhd}               & \texttt{alu / rtl}                 & 0 & 0 \\
    \texttt{fase2/timer.vhd}             & \texttt{timer / rtl}               & 0 & 0 \\
    \texttt{fase2/clk\_div.vhd}          & \texttt{clk\_div / rtl}            & 0 & 0 \\
    \texttt{fase2/ctrl\_tec.vhd}         & \texttt{ctrl\_tec / rtl}           & 0 & 0 \\
    \texttt{fase2/interfaz\_teclado.vhd} & \texttt{interfaz\_teclado / est.}  & 0 & 0 \\
    \texttt{fase2/displays.vhd}          & \texttt{displays / rtl}            & 0 & 0 \\
    \texttt{fase2/tb\_alu.vhd}           & \texttt{tb\_alu / sim}             & 0 & 0 \\
    \texttt{fase3/controlador.vhd}       & \texttt{controlador / rtl}         & 0 & 0 \\
    \texttt{fase3/calculadora.vhd}       & \texttt{calculadora / estructural} & 0 & 0 \\
    \texttt{fase3/tb\_calculadora.vhd}   & \texttt{tb\_calculadora / sim}     & 0 & 0 \\
    \midrule
    \textbf{Total}                       & & \textbf{0} & \textbf{0} \\
    \bottomrule
  \end{tabular}
\end{table}

La simulación del testbench \texttt{tb\_alu} finaliza con el mensaje:

\begin{verbatim}
  # ** Note: tb_alu: todos los casos completados
  # ** Failure: Assertion failure
\end{verbatim}

El \textit{Assertion failure} al final es el mecanismo de parada controlada
(\texttt{assert false severity failure}), no un error del diseño.
\FloatBarrier

% ===================================================================
\section{Conclusiones}
% ===================================================================

En esta fase se han diseñado, implementado y verificado todos los bloques
intermedios de la calculadora:

\begin{itemize}[leftmargin=1.6cm,itemsep=5pt]
  \item \textbf{ALU} (\texttt{alu.vhd}): lógica combinacional pura,
    implementada íntegramente con asignaciones concurrentes
    \texttt{when/else}. Maneja suma, resta y multiplicación con signo en
    complemento a~2, extrayendo la magnitud del resultado sin variables.

  \item \textbf{Temporizadores} (\texttt{clk\_div}, \texttt{timer}):
    generan tics de 5\,ms para el teclado y de 1\,ms para el mux de
    displays, con genéricos que permiten reducir los valores en simulación.

  \item \textbf{Controlador de teclado} (\texttt{ctrl\_tec}): antirebote
    integrado de dos etapas (5\,ms de estabilidad), escaneo one-hot de filas,
    detección de pulsación en el flanco de bajada, soporte de pulsación larga
    (2\,s) con supresión del evento de release.

  \item \textbf{Controlador de displays} (\texttt{displays}): multiplexación
    a 125\,Hz, supresión de ceros a la izquierda con asignaciones concurrentes,
    soporte de signo negativo.
\end{itemize}

El criterio de no usar variables se ha aplicado en todos los módulos RTL.
Toda la lógica combinacional es visible como asignaciones concurrentes fuera
de los procesos, lo que facilita la revisión del diseño y la correspondencia
directa con el hardware sintetizado.

\vspace{1em}
\begin{tcolorbox}[colback=SoftBlue,colframe=MainBlue,boxrule=0.8pt,arc=2pt,
  left=0.5cm,right=0.5cm,top=0.3cm,bottom=0.3cm]
  \small\faArrowRight\enspace\textbf{Fase 3 (entrega 28/05):}
  Diseño del módulo \texttt{controlador} (FSM principal) y del top-level
  \texttt{calculadora}, que integra todos los bloques de las tres fases en
  un sistema completo verificable con \texttt{tb\_calculadora}.
\end{tcolorbox}

% ===================================================================
\newpage
\appendix
\section{Señales frente a variables en VHDL}
\label{sec:senales-variables}
% ===================================================================

\subsection{Diferencia semántica}

En VHDL existen dos mecanismos de almacenamiento temporal dentro de un proceso.

\paragraph{Señal (\texttt{signal}):} La asignación \texttt{s <= expr} no
produce ningún efecto inmediato. El nuevo valor de \texttt{s} solo es visible
\textit{después} de que el proceso termine (en el siguiente delta-cycle).
Esto modela directamente el comportamiento de un registro físico: la salida
de un flip-flop solo cambia en el flanco de reloj, no en medio del ciclo.

\paragraph{Variable (\texttt{variable}):} La asignación \texttt{v := expr}
actualiza \texttt{v} inmediatamente. El siguiente uso de \texttt{v} en ese
mismo proceso ya ve el valor nuevo. Este comportamiento es propio del
software; en hardware no tiene equivalente directo, y el sintetizador debe
inferir qué circuito produce ese comportamiento.

\subsection{Ejemplo comparativo}

\begin{lstlisting}[caption={Shift register con señales (correcto)}]
process(clk)
begin
    if clk'event and clk = '1' then
        a <= '1';   -- a todavia vale '0' aqui
        b <= a;     -- b recibe el valor ANTIGUO de a
        c <= b;     -- c recibe el valor ANTIGUO de b
    end if;         -- al terminar: a='1', b=a_viejo, c=b_viejo
end process;        -- -> shift register de 3 etapas (correcto)
\end{lstlisting}

\begin{lstlisting}[caption={Mismo código con variables (resultado diferente)}]
process(clk)
    variable a, b, c : std_logic;
begin
    if clk'event and clk = '1' then
        a := '1';   -- a vale '1' AHORA MISMO
        b := a;     -- b recibe '1' (valor nuevo)
        c := b;     -- c recibe '1' tambien
    end if;         -- resultado: a=b=c='1' en el mismo ciclo
end process;        -- -> NO es un shift register
\end{lstlisting}

\subsection{Consecuencias prácticas}

\begin{table}[H]
  \centering
  \caption{Comparación señales vs.\ variables en diseño RTL}
  \begin{tabular}{lll}
    \toprule
    \textbf{Aspecto} & \textbf{Señal} & \textbf{Variable} \\
    \midrule
    Actualización      & Al terminar el proceso    & Inmediata \\
    Modelado HW        & Registro / cable           & Memoria temporal software \\
    Correspondencia HW & Directa y explícita        & Inferida por el sintetizador \\
    Visible en waveform& \textbf{Sí}               & \textbf{No} \\
    Depuración         & Sencilla                  & Difícil \\
    Estilo RTL         & Recomendado IEEE           & Desaconsejado en síncronos \\
    \bottomrule
  \end{tabular}
\end{table}

\subsection{Regla práctica adoptada}

En este proyecto se sigue la siguiente convención de forma estricta:

\begin{enumerate}[leftmargin=1.8cm,itemsep=3pt]
  \item \textbf{Dentro del proceso}: solo asignaciones \texttt{<=} a
    registros (flip-flops). La condición de activación siempre es
    \texttt{clk'event and clk = '1'}.
  \item \textbf{Fuera del proceso}: toda la lógica combinacional como
    asignaciones concurrentes con \texttt{<=} o \texttt{when/else}.
\end{enumerate}

Esta separación hace que el código sea autoexplicativo: leer los procesos
muestra los registros del sistema, y leer las asignaciones concurrentes
muestra la lógica combinacional entre ellos.

\end{document}
